% Breakable frames preserve prompt contents while using the available page space.
\definecolor{promptframeborder}{gray}{0.55}
\definecolor{promptframebackground}{gray}{0.975}
\newenvironment{promptframe}[1]{%
  \begin{tcolorbox}[enhanced,breakable,
    colback=promptframebackground,colframe=promptframeborder,
    boxrule=0.5pt,arc=0pt,left=9pt,right=9pt,top=6pt,bottom=6pt,
    before skip=8pt,after skip=8pt,
    title={#1},title after break={#1 (continued)},
    colbacktitle=gray!10,coltitle=black,fonttitle=\small\bfseries,
    fontupper=\small,lines before break=4]
  \raggedright\setstretch{1.0}\setlength{\parindent}{0pt}\setlength{\parskip}{4pt}%
}{%
  \end{tcolorbox}%
}

\section{Prompts}
\label{app:prompts}
This section gives compact prompt templates and one recorded interaction for
each benchmark. Templates summarize the implemented instructions; bracketed
fields are populated from the task package and run configuration. They are
not verbatim copies of the complete runtime prompts, which also contain
structured context, tool definitions, and conversation history. The simulator
plays the researcher, and only its prompt receives the persona profile.
The examples paraphrase and abridge the shared messages while retaining every
researcher review round or clarification interaction through termination (at
most three per example). Tool calls and private evaluator messages are
omitted. Numerical values in the examples are actor-reported results, not
independent validations or evidence that the submission received full credit.

For DiscoveryBench, BixBench, EXP-Bench, and ASI-Bench, each actor handoff
contains its stage (plan or result), answer, evidence, assumptions, and
uncertainties, together with any required artifacts. The simulator returns
a continue/stop decision, a message, the remaining objectives, any disclosed
fact identifiers, and an approach-approval flag. It can stop only during
result review. SciConvBench instead retains its native question-and-answer
interface. These interfaces implement the boundaries described in
Appendix~\ref{app:conversion-principles}.

\subsection{Example Persona Profile: Persona 155}
\label{app:persona-profile}
Persona 155 (\texttt{candidate\_155}) is one of the four injected profiles
used in the current analysis (Section~\ref{sec:run-coverage}). The frames
below reproduce all fields of its configured identity block, preserving
their wording, values, and order while replacing the original line-by-line
layout with grouped text. The profile is supplied only to the simulator;
the task context and common protocol instructions are supplied separately.
These are synthetic profile attributes, not verified biographical facts or
additional scientific task requirements.

% Exact source: configs/main_run/personas/candidate_155.md
% Matching rendered profile: reports/full_persona_v2/profiles/candidate_155.md
% Source SHA256: 8612aa84a205eb0de938f82b4a7bd4537cdd94f4e93bef4f30b4d694835a7652
\begin{promptframe}{Persona 155: Who you are}
\textbf{Identity}\par\nobreak
Age bracket: 85+; Region: Eastern Europe; Gender identity: Woman; Urbanicity: Rural; Socioeconomic band: Upper-middle; Dialect / register: Formal / standard; Life stage: Retirement; Cultural background: Collectivist (East Asian); Household income band: \$100k-200k; Employment status: Retired; Housing status: Own outright; Generational cohort: Silent; Citizenship status: Citizen by birth; Ethnic background: White / European; Disability status: No disability; Veteran status: Civilian; Language at home: Same as primary; Political engagement: Activist; Driving status: Non-driver; Childhood environment: Modest but steady; Geographic mobility: Moved within region; Travel breadth: A few countries; Education journey: Vocational / trade; Adversity faced: Minor setbacks; Immigration generation: Native multi-generational; Faith journey: Always secular; Formative decade: 1920s; Public recognition: Industry-recognized; Community / service history: Long-term volunteer; Defining turning point: No single one\par
\textbf{Career \& education}\par\nobreak
Highest education: No formal; Seniority: Manager; Years experience: 20+\par
\textbf{Language \& communication}\par\nobreak
Primary language: Mandarin; Multilingualism: Monolingual; Detail orientation: High; Abstract vs concrete: Concrete; Numeracy comfort: High; Emoji / emoticon use: Never; Default written language: use your primary language for outputs.\par
\textbf{Personality \& values}\par\nobreak
Domain stance: Practitioner; Dominant trait: High conscientiousness; Core value: Achievement; Religiosity: Secular; Economic motivation: Value-driven; Self-efficacy: High; Orderliness: High; Dutifulness: Very high; Achievement-striving: Very high; Self-discipline: High; Assertiveness: High; Activity level: High; Cooperation: High; Character: Perseverance: Signature; Character: Zest: Strong; Character: Teamwork: Signature; Character: Leadership: Signature; Character: Prudence: Strong; Character: Self-regulation: Signature; Character: Ambition: Signature; Character: Resilience: Signature; Character: Discipline: Signature; Character: Loyalty: Signature; Character: Competitiveness: Strong; Character: Adaptability: Strong; Value: Career success: Core value; Value: Security \& stability: Important; Value: Tradition: Important; Value: Achievement: Core value; Value: Community: Important; Value: Social status: Important; Value: Loyalty: Important; Value: Recognition: Core value; Value: Order \& structure: Important; Value: Patriotism: Core value; BFI-2 Conscientiousness: High; BFI-2 Assertiveness: High; BFI-2 Energy Level: High; BFI-2 Organization: High; BFI-2 Productiveness: Very high; BFI-2 Responsibility: Very high; Schwartz Self-Direction: Low; Schwartz Stimulation: Low; Schwartz Hedonism: Low; Schwartz Achievement: Very high; Schwartz Security: High; Schwartz Conformity: High; Schwartz Universalism: Low; SDT Autonomy Need: Low; SDT Competence Need: Very high; SDT Relatedness Need: High; Need for Cognition: Low; Interpersonal Agency/Dominance: High\par
\textbf{Current interaction state}\par\nobreak
Expertise gap: Expert testing the system; Trust level: Trusting; Safety sensitivity: Benign; Work schedule: 9-to-5; Office vs remote: Strongly office\par
\textbf{Worldview}\par\nobreak
Political lean: Left; Familiarity: Political science: Familiar; Attitude: Organized religion: Skeptical; Attitude: Labor unions: Positive; Attitude: Authority: Positive\par
\textbf{Interests}\par\nobreak
Interest: Politics: Interested; Interest: Travel: Interested; Interest: Activism: Interested; Interest: Productivity: Passionate; Culture: Russia: Lived there\par
\textbf{Skills \& expertise}\par\nobreak
Domain: Agriculture; Subject specialty: Agronomy; Tech savviness: Reluctant; Familiarity: Statistics: Familiar; Familiarity: Accounting: Familiar; Familiarity: Agronomy: Proficient; Familiarity: Horticulture: Familiar; Familiarity: Supply chain: Aware; Familiarity: Logistics: Aware; Familiarity: Operations management: Familiar; Familiarity: Quality assurance: Aware; Familiarity: Human resources: Aware; Familiarity: Project management: Aware; Familiarity: Textiles: Aware; Familiarity: Meteorology: Aware; Familiarity: Materials science: Aware; Skill: Project management: Intermediate; Skill: Leadership: Advanced; Skill: People management: Advanced\par
\textbf{Lifestyle \& health}\par\nobreak
Team vs solo work: Strongly team; City vs nature: Prefers nature; Speed vs accuracy: Accuracy first; Quality vs quantity: Quality first; Novelty vs familiarity: Always familiar; Competition vs collaboration: Highly collaborative; Stability vs change: Craves stability\par
\textbf{Developer \& AI}\par\nobreak
Developer professional status: Non-developer technologist\par
\end{promptframe}

\subsection{DiscoveryBench}
\label{app:prompts-discoverybench}
\begin{promptframe}{Prompt template}
\textbf{Actor:} Analyze [question] using [public datasets and column
metadata]. Inspect the inputs and hand off a proposed analysis with its
assumptions before proceeding. After researcher approval, perform the
analysis and report a data-grounded hypothesis, evidence, and limitations.
Preserve the original question and fixed constraints.\par
\textbf{Simulator:} You are the researcher described by [persona, if enabled].
Use [public brief], [approved domain-knowledge facts], and [shared history]
as task context. Request inspection first, review the proposed choices,
and approve or seek a focused clarification. Review the reported hypothesis
against the agreed objective; request a specific revision or conclude.
Distinguish methodological advice from established study facts. Do not
invent reference answers or grade the hypothesis.
\end{promptframe}

% Source: reports/full_persona_v2/recorded_traces/candidate_055-5636a4f47387.md
\begin{promptframe}{Recorded example: symbolic capital and PCA}
Task: DiscoveryBench archaeology, metadata 21, question 0;
trace: \texttt{candidate\_055-5636a4f47387}.
\par\smallskip
\textbf{Round 1---Simulator:} Inspect the capital data and any existing PCA
code. Propose an approach for house size, daggers, and hatchets/swords;
identify uncertainties without computing the full result.\par
\textbf{Actor:} The table has 39 hundred-year slices and 12 capital indicators,
with decimal-comma values and no existing PCA implementation. I propose
standardizing the indicators and reporting loadings separately from
observation scores. The original scaling and component signs are uncertain.\par
\textbf{Round 2---Simulator:} Proceed with standardized PCA. State the sign
convention and identify this as a reconstruction; optionally check the
unstandardized result.\par
\textbf{Actor:} The reported PC1 loadings are 0.429, 0.507, and 0.570 for house
size, daggers, and hatchets/swords. PC2 loadings are 0.500, 0.113, and 0.082.
The indicators align on PC1, whereas PC2 distinguishes house size from the
weapon indicators. The original PCA specification remains undocumented.\par
\textbf{Round 3---Simulator:} Conclude with that hypothesis, retaining the
reconstruction caveat and the distinction between loadings and scores.
\end{promptframe}

\subsection{BixBench}
\label{app:prompts-bixbench}
\begin{promptframe}{Prompt template}
\textbf{Actor:} Answer [original question] using [capsule inputs]. Inspect
the schema and propose the analysis before requesting approval. Preserve
specified methods and thresholds. After execution, hand off the answer and
evidence, provide [analysis script or notebook], and set the execution
declaration to true only if the analysis was run.\par
\textbf{Simulator:} Use [persona, if enabled], [public question], [approved
input-derived metadata], and [shared history]. Review consequential choices
such as column mapping and contrast direction. Treat undocumented meanings
as unknown and proposed analysis choices as assumptions. Approve the
approach when sufficiently specified, then review the reported answer and
required execution artifact before concluding or requesting a revision.
\end{promptframe}

% Source: reports/full_persona_v2/recorded_traces/candidate_138-dedb31fc5cc4.md
\begin{promptframe}{Recorded example: differential expression by tissue}
Task: \texttt{bix-9-q5}; trace: \texttt{candidate\_138-dedb31fc5cc4}.
\par\smallskip
\textbf{Round 1---Simulator:} Inspect response groups, tissue labels, and
expression-result fields. Propose how to compare bad responders with
controls using the specified fold-change and mean-expression thresholds.\par
\textbf{Actor:} The workbook contains sample annotations and precomputed
DESeq results. I propose counting genes within each tissue using the signed
\texttt{log2FoldChange} $>0.5$ and \texttt{baseMean} $>10$ thresholds.
The significance cutoff is unspecified; I propose adjusted $p<0.05$, subject
to checking the workbook's definitions.\par
\textbf{Round 2---Simulator:} Proceed after verifying the comparison labels
and contrast direction. Use adjusted $p<0.05$ unless the workbook documents
another criterion, state the definition, and retain the executed script.\par
\textbf{Actor:} Dentate gyrus has 194 qualifying genes, compared with 88 in
baseline blood and 14 in final blood, using the verified comparison blocks
and those thresholds.\par
\textbf{Round 3---Simulator:} Conclude that dentate gyrus has the highest
count. Retain the thresholds and the executed \texttt{analysis/analyze.py}
as evidence.
\end{promptframe}

\subsection{EXP-Bench}
\label{app:prompts-expbench}
\begin{promptframe}{Prompt template}
\textbf{Actor:} Develop an experimental design for [scientific question]
using [public task material]. Inspect what is actually available, propose
an approach, and wait for approval. Then present the complete design with
rationale, controls, analysis, feasibility, assumptions, and limitations.
This is design-only: do not claim experimental execution or invent outcomes.\par
\textbf{Simulator:} Use [persona, if enabled], [public brief], [approved
setup facts], and [shared history]. Review the approach before the full
design. Distinguish fixed setup requirements from choices the actor must
make. Review the completed design for alignment and unresolved uncertainties;
request a specific revision or conclude without demanding executed results.
\end{promptframe}

% Source: reports/full_persona_v2/recorded_traces/candidate_224-e1ecc5b8f19c.md
\begin{promptframe}{Recorded example: comparing MAPPO across VMAS tasks}
Task: \texttt{neurips2024/98318/6};
trace: \texttt{candidate\_224-e1ecc5b8f19c}.
\par\smallskip
\textbf{Round 1---Simulator:} Inspect the available material and propose a
design comparing MAPPO on VMAS balance and sampling. Identify assumptions;
do not produce a full design or infer an empirical outcome yet.\par
\textbf{Actor:} Only the task description is present, with no repository or
task configuration. I propose matching training budgets and evaluation
protocols, reporting task-specific performance and uncertainty, and checking
reward semantics before comparing raw rewards across tasks.\par
\textbf{Round 2---Simulator:} Proceed with the full design. Separate documented
facts from provisional choices, and cover metrics, controls, feasibility,
and unresolved implementation details without claiming execution.\par
\textbf{Actor:} First establish the intended versions and configurations.
Use the same MAPPO implementation and equal interaction budgets, repeated
training runs, and matched held-out evaluation. Report task-specific success
or returns, learning curves, uncertainty, and failures. Qualify cross-task
reward comparisons; numerical settings remain contingent on the missing
specification and resources.\par
\textbf{Round 3---Simulator:} The design adequately addresses the comparison,
controls, reward-scale differences, and remaining uncertainties. Conclude
without claiming an executed experiment or expected outcome.
\end{promptframe}

\subsection{SciConvBench}
\label{app:prompts-sciconvbench}
\begin{promptframe}{Prompt template}
\textbf{Actor:} Clarify [incomplete or inconsistent scientific request]. Ask
one focused question at a time and wait for the answer. Address missing or
contradictory requirements. When ready, return \texttt{[COMPLETE]} followed
by the final specification as plain sentences.\par
\textbf{Simulator:} You know [complete or corrected requirement], whereas
the actor sees [initial request]. Answer clarification questions using the
intended requirement and the configured knowledge policy. Give concrete
values when supplied; do not turn unspecified details into source facts.
[Persona, if enabled] sets how you communicate and does not change the
requirements. [Shared dialogue] supplies the conversation history.
\end{promptframe}

% Source: reports/full_persona_v2/recorded_traces/candidate_046-b5331aa3b9ba.md
\begin{promptframe}{Recorded example: elongation under self-weight}
Task: solid mechanics, disambiguation case 104;
trace: \texttt{candidate\_046-b5331aa3b9ba}.
\par\smallskip
\textbf{Initial request:} Determine the elongation of a vertical bar under
self-weight, assuming small-deformation linear elasticity, and state the
governing relations and any additional assumptions.\par
\textbf{Interaction 1---Actor:} Should the bar be uniform and prismatic, fixed
at its top, free at its bottom, and loaded only by its own weight?\par
\textbf{Simulator:} Yes. It hangs from a rigid top support with a free bottom
and no external load. At distance $x$ downward from the top,
$P(x)=\rho gA(L-x)$.\par
\textbf{Actor---final specification:} Use a uniform bar of length $L$, area
$A$, density $\rho$, and Young's modulus $E$, with constant properties and
gravity $g$. Fix the top and leave the bottom free. With $u(0)=0$, use
$\sigma(x)=\rho g(L-x)$ and
$u(x)=(\rho g/E)(Lx-x^2/2)$, giving elongation
$\delta=\rho gL^2/(2E)$. Use consistent units.
\par
The actor terminated after one clarification interaction. This records its
submitted specification, not a claim that all intended requirements were
recovered; no solver execution accompanies the dialogue.
\end{promptframe}

\subsection{ASI-Bench}
\label{app:prompts-asibench}
\begin{promptframe}{Prompt template}
\textbf{Actor:} Complete [native B3 task] using [permitted inputs] and produce
[required output files]. Inspect the allowed inputs, propose an approach,
and wait for approval before implementation. Execute the approved work and
hand off the results, diagnostics, limitations, and artifact paths. Preserve
all task constraints and access restrictions.\par
\textbf{Simulator:} Use [persona, if enabled], [B3 public brief], [approved
B2 methodological guidance], and [shared history]. Treat guidance as advice
rather than a new fixed requirement. Review the proposed method and then
the reported outputs against the original contract. Request a specific
revision or conclude; do not claim access to private reference results or
treat actor reports as independently verified evidence.
\end{promptframe}

% Source: reports/full_persona_v2/recorded_traces/candidate_137-bdb8b6dd4a8c.md
\begin{promptframe}{Recorded example: nonnegative deblurring}
Task: \texttt{math.nnls\_modulus\_deblur};
trace: \texttt{candidate\_137-bdb8b6dd4a8c}.
\par\smallskip
\textbf{Round 1---Simulator:} Inspect only permitted inputs. Report their
structure and propose an iterative nonnegative reconstruction with
regularization, diagnostics, and the required artifacts; do not implement yet.\par
\textbf{Actor:} The observation and kernel are $128\times128$ arrays.
I propose FFT forward and adjoint operators with projected accelerated
gradient, a regularization sweep, and convergence diagnostics. Plotting
libraries are absent, and the noise estimate is only approximate.\par
\textbf{Round 2---Simulator:} Proceed with the implementation and validate
all five required outputs. Preserve the supplied kernel, document the
regularization and stopping rule, and resolve PNG generation within the
permitted environment.\par
\textbf{Actor:} The reconstruction ran with $\lambda=10^{-4}$ and stopped
after 1,335 iterations. The reported residual is 2.16684884 against an
approximate target of 2.358912. The script, reconstruction array, PNG,
iteration log, and metrics file are present and pass the reported format
checks. No ground-truth error is claimed.\par
\textbf{Round 3---Simulator:} Conclude on the basis of the reported artifacts
and diagnostics, retaining the approximate-noise and ground-truth limitations.
\end{promptframe}
